% Event-study subsection for manual inclusion in paper/main.tex
% Source tables:
% - out_files/ftsemib/metaorder_start_event_study/event_study_summary.csv
% - out_files/ftsemib/metaorder_start_event_study/event_study_curves.csv
% - out_files/ftsemib/metaorder_start_event_study/event_study_diagnostics.csv
% - out_files/ftsemib/metaorder_start_event_study/robustness_same_actor_exclusion_summary.csv

\subsection{Local Start-Time Crowding Around High-Participation Metaorders}
\label{subsec:start_event_study}

The day-level crowding statistics of Section~\ref{subsec:imbalance} show that
metaorders tend to align with the contemporaneous imbalance generated by other
orders in the same group. A complementary question is whether this crowding is
already visible in the local neighbourhood of the \emph{start} of a metaorder,
and whether it is especially pronounced for the most aggressive executions. To
address this point, we run an event-study centred on the start time
$t_i^{\mathrm{start}}$ of each metaorder $i$.

For each group $G \in \{\text{prop},\text{client}\}$, we classify as
high-participation anchors the metaorders with participation rate $\eta_i$
above the group-specific median. In the present sample this threshold equals
$\eta_{\mathrm{prop}}^{\star}=0.09499$ for proprietary flow and
$\eta_{\mathrm{client}}^{\star}=0.05019$ for client flow. Around each anchor
$i$, we collect all other metaorder starts $j$ in the same ISIN and on the same
day and define the event time
\begin{equation}
  \tau_{ij} = t_j^{\mathrm{start}} - t_i^{\mathrm{start}}.
\end{equation}
We then measure, separately for same-sign ($\varepsilon_j=\varepsilon_i$) and
opposite-sign ($\varepsilon_j=-\varepsilon_i$) starts, the start intensity in
5-minute bins over the window $\tau \in [-30,30]$ minutes. To remove slow
intraday seasonality, each high-$\eta$ anchor is compared with matched controls
chosen within the same ISIN, date, and 30-minute clock-time bucket. Event-time
rates are normalised by the effective exposure time in each bin, so that
session-boundary truncation does not mechanically bias the intensity estimates.

The baseline results point to a clear post-start asymmetry. For proprietary
anchors, the mean same-sign start intensity over the post-event window exceeds
that of matched controls by $0.0014$ starts per minute, with Holm-adjusted
$p=0.004$, while the opposite-sign intensity is lower by $0.0015$ starts per
minute, with Benjamini--Hochberg-adjusted $p=0.002$. For client anchors, the
same pattern is considerably stronger: the post-event same-sign excess is
$0.0076$ starts per minute (Holm-adjusted $p=0.004$), whereas the post-event
opposite-sign intensity is lower by $0.0082$ starts per minute
(Benjamini--Hochberg-adjusted $p=0.002$). Hence, high-participation starts are
followed by a more one-sided local order-flow environment in both segments, but
the effect is materially larger for client flow. The corresponding event-time
curves are shown in Figure~\ref{fig:start_event_study_baseline}.

The event-time profiles reveal that this effect is not concentrated in the
first few minutes after the anchor. In both groups, the $(0,5]$ minute bin
displays a deficit of same-sign starts relative to matched controls, but this
is followed by a persistent excess from 5 to 30 minutes after the start.
Opposite-sign starts, instead, fall sharply in the first post-event bin and
remain below control levels over most of the subsequent half-hour. On the
pre-event side, the client curves already display elevated same-sign intensity
and depressed opposite-sign intensity over much of the interval $[-30,-5)$,
suggesting that high-$\eta$ client metaorders tend to start in already one-sided
local environments. The analogous proprietary pre-pattern is weaker and does
not survive multiple-testing adjustment at conventional levels.

\begin{figure}[t]
  \centering
  \begin{adjustwidth}{-1.8cm}{-1.8cm}
    \centering
    \begin{subfigure}[t]{0.60\textwidth}
      \centering
      \includegraphics[width=\textwidth]{images/metaorder_start_event_study/png/event_curve_prop_all_others.png}
      \caption{Proprietary flow.}
    \end{subfigure}
    \hfill
    \begin{subfigure}[t]{0.60\textwidth}
      \centering
      \includegraphics[width=\textwidth]{images/metaorder_start_event_study/png/event_curve_client_all_others.png}
      \caption{Client flow.}
    \end{subfigure}
  \end{adjustwidth}
  \caption{Matched start-time event-study around high-participation metaorders.
  Each panel reports the start intensity of neighbouring same-sign and
  opposite-sign metaorders in event time, comparing high-$\eta$ anchors with
  matched controls from the same ISIN, date, and 30-minute clock-time bucket.
  The event window is $\pm 30$ minutes in 5-minute bins.}
  \label{fig:start_event_study_baseline}
\end{figure}

A key robustness check excludes neighbouring starts generated by the same
actor. At the member level, the relevant actor identifier is the executing
member, so the exclusion removes starts initiated by the same member as the
anchor. Under this stricter specification, the same-sign excess largely
disappears and, if anything, turns negative for both proprietary and client
flow. By contrast, the post-event deficit of opposite-sign starts remains
strong: $-0.0082$ starts per minute for proprietary anchors and $-0.0107$
starts per minute for client anchors, with permutation $p$-values close to
$10^{-3}$ in both cases. This robustness check changes the interpretation of
the baseline event study. The raw same-sign clustering around
high-participation starts is driven to a large extent by repeated start
activity of the same member, rather than by broad cross-member synchronisation.
What remains after excluding same-member neighbours is mainly a depletion of
opposite-side starts around the anchor. Figure~\ref{fig:start_event_study_same_actor}
shows that this robustness adjustment materially flattens the same-sign curves
for both groups while leaving the post-event deficit of opposite-sign starts
clearly visible.

\begin{figure}[t]
  \centering
  \begin{adjustwidth}{-1.8cm}{-1.8cm}
    \centering
    \begin{subfigure}[t]{0.60\textwidth}
      \centering
      \includegraphics[width=\textwidth]{images/metaorder_start_event_study/png/event_curve_prop_exclude_same_actor.png}
      \caption{Proprietary flow, excluding starts from the same member.}
    \end{subfigure}
    \hfill
    \begin{subfigure}[t]{0.60\textwidth}
      \centering
      \includegraphics[width=\textwidth]{images/metaorder_start_event_study/png/event_curve_client_exclude_same_actor.png}
      \caption{Client flow, excluding starts from the same member.}
    \end{subfigure}
  \end{adjustwidth}
  \caption{Robustness event-study excluding neighbouring starts generated by
  the same executing member as the anchor. Relative to
  Figure~\ref{fig:start_event_study_baseline}, the excess of same-sign starts
  is strongly attenuated, whereas the short-horizon depletion of opposite-sign
  starts remains.}
  \label{fig:start_event_study_same_actor}
\end{figure}

The sample sizes are large, but the matching design is non-trivial. The event
study uses $588{,}334$ proprietary and $255{,}535$ client anchors overall;
among them, $294{,}167$ proprietary and $127{,}768$ client anchors are above
the high-$\eta$ threshold. Exact matching retains $230{,}342$ proprietary and
$79{,}633$ client treated anchors, implying that about $21.7\%$ of proprietary
and $37.7\%$ of client high-participation anchors are dropped because no valid
control is available in the same $(\mathrm{ISIN}, \mathrm{Date},
\mathrm{bucket})$ stratum. Confidence intervals are obtained by a day-cluster
bootstrap and $p$-values by within-stratum permutation, so the evidence should
be read as descriptive but dependence-aware rather than as a causal estimate of
the effect of a metaorder start.

Taken together, these results complement the day-level crowding analysis.
High-participation metaorders, especially client ones, start in local
environments that are already directionally unbalanced and are followed by a
short-horizon continuation in the same direction together with a contraction of
opposite-sign starting activity. However, once same-member repetition is
removed, the continuation result weakens substantially. The most robust local
signature of a high-participation start is therefore not generalised same-side
herding, but a temporary scarcity of opposite-side starts in its immediate
neighbourhood.
